\documentclass[11pt,a4paper]{article}
\usepackage[T1]{fontenc}
\usepackage[utf8]{inputenc}
\usepackage{amsmath,amssymb,xcolor}
\usepackage[margin=25mm]{geometry}
\usepackage[colorlinks=true,urlcolor=blue]{hyperref}
\definecolor{nexorablue}{HTML}{2459B5}
\setlength{\parindent}{0pt}
\setlength{\parskip}{7pt}
\setlength{\emergencystretch}{2em}
\renewcommand{\familydefault}{\sfdefault}
\begin{document}

{\small\bfseries\color{nexorablue} NEXORA CPGE}\par
{\small MP / Cours}\par
{\LARGE\bfseries Modélisation des systèmes asservis}\par
\vspace{0.25em}\hrule\vspace{0.6em}

\section*{Objectifs et prérequis}
Revoir les équations différentielles linéaires et les fractions rationnelles. On cherche un modèle entrées-sorties d’un système linéaire, continu et invariant autour d’un point de fonctionnement. Distinguer la consigne, l’action de commande, la sortie mesurée et les perturbations.

\section*{Modèle et fonction de transfert}
Pour un système scalaire, prendre la transformée de Laplace unilatérale de l’équation qui relie les signaux d’entrée $u$ et de sortie $y$. La fonction de transfert correspond à des conditions initiales nulles, lorsqu’on décrit les écarts au point de fonctionnement.
\[ H(p)=\frac{Y(p)}{U(p)},\qquad \mathcal L\{\dot y\}=pY(p)-y(0^-). \]
Par exemple, $\tau\dot y+y=Ku$ conduit à $H=K/(1+\tau p)$ sous ces hypothèses. Avec $y(0^-)\ne0$, traiter séparément la contribution de l’état initial.

\section*{Combiner des blocs}
En série, la sortie du premier bloc devient l’entrée du second. On multiplie les transferts. En parallèle, si deux branches reçoivent la même entrée, le sommateur ajoute ou retranche leurs sorties selon les signes indiqués.
\[ \text{série : }H=H_2H_1,\qquad \text{parallèle signé : }H=H_1\pm H_2. \]
Un prélèvement copie la grandeur à un nœud. Avant de déplacer un sommateur ou un prélèvement, vérifier que les deux schémas conservent exactement les mêmes équations.

\section*{Retour négatif}
Soit $R$ la consigne exprimée dans l’unité du signal mesuré, $G$ la chaîne directe et $B$ le transfert du capteur. Le comparateur forme $E=R-BY$. La chaîne directe donne $Y=GE$. Substituer la première équation dans la seconde, puis isoler la sortie.
\[ \frac{Y}{R}=\frac{G}{1+GB},\qquad \frac{E}{R}=\frac1{1+GB}. \]
Le produit $GB$ est le transfert de boucle. En retour positif, le même raisonnement donne un dénominateur $1-GB$. Pour un retour unitaire, prendre $B=1$.

\section*{Perturbation et stabilité}
Si $D$ s’ajoute en sortie du procédé, écrire $Y=G(R-BY)+D$. On obtient $Y_D=D/(1+GB)$ à consigne nulle. Identifier l’emplacement de la perturbation avant de choisir la formule.

En temps continu, les pôles d’une fonction rationnelle simplifiée doivent avoir une partie réelle strictement négative pour assurer la stabilité entrée-sortie. Une simplification de facteurs peut masquer un mode interne. Vérifier la stabilité avant d’appliquer le théorème de la valeur finale.

\section*{Réponse à un échelon et régime permanent}
Une consigne en échelon d’amplitude $R_0$ a pour transformée $R_0/p$. Calculer $Y(p)$, puis effectuer une décomposition en éléments simples. Si les pôles pertinents de $pY(p)$ sont strictement à gauche, $y(\infty)=\lim_{p\to0}pY(p)$.

Pour $H(p)=K/(1+\tau p)$ avec $\tau>0$, la sortie initialement nulle vaut $KR_0(1-e^{-t/\tau})$. Le temps de réponse à $5\,\%$ se lit en résolvant $e^{-t/\tau}=0{,}05$, soit environ $3\tau$.

\section*{Exemple et réflexes CNC}
Boucler $G=2/(1+p)$ avec un capteur unitaire donne $H=2/(p+3)$. La constante de temps devient $1/3$ et le gain statique $2/3$. Pour un échelon unitaire, l’écart permanent vaut $1/3$. Une réponse plus rapide n’implique donc pas une erreur statique nulle.

Sur une copie, nommer les unités des deux signaux comparés, écrire les équations du comparateur et de la sortie, puis identifier le dénominateur avant de discuter les pôles.

\section*{Passer d’un modèle physique au schéma-blocs}
Identifier d’abord les entrées indépendantes, par exemple une tension de commande et un couple résistant. Écrire une équation pour chaque relation électrique et mécanique, puis prendre la transformée sous conditions initiales nulles. Construire les blocs à partir de ces équations, sans perdre les signes.

Pour un moteur à courant continu sans inductance significative, $u=Ri+k_e\omega$ et $J\dot\omega+f\omega=k_t i-C_r$. La vitesse influe sur le courant via la force contre-électromotrice, même avant l’ajout d’une boucle de régulation externe.

\section*{Distinguer consigne et perturbation}
Un système peut avoir plusieurs entrées et une seule sortie. Pour trouver un transfert, annuler les autres entrées sans supprimer physiquement les blocs. Une perturbation $D$ ajoutée à la sortie d’une chaîne directe $G$ donne, avec le capteur $B$,
\[ Y=\frac{G}{1+GB}R+\frac{1}{1+GB}D. \]
Si $D$ est additive à l’entrée de $G$, le second terme devient $GD/(1+GB)$. Ces deux modèles n’ont pas les mêmes unités pour $D$. Nommer l’endroit où elle apparaît sur le schéma.

\section*{Régler un retour sur un modèle du premier ordre}
Avec $G(p)=K/(1+\tau p)$ et un capteur constant $b$, la boucle négative a pour transfert
\[ H(p)=\frac{K}{1+Kb+\tau p}. \]
Pour $\tau>0$ et $1+Kb>0$, on lit un gain statique $K/(1+Kb)$ et une constante de temps $\tau/(1+Kb)$. L’erreur en régime permanent porte sur $R-bY$, c’est-à-dire sur deux signaux exprimés dans les mêmes unités.

Comparer chaque exigence du cahier des charges séparément. Une contrainte sur la rapidité et une autre sur la précision produisent souvent deux bornes différentes sur le gain.

\section*{Contrôles avant de conclure}
Revenir aux limites : si le gain direct tend vers zéro, le suivi de consigne doit disparaître. Pour un système stable, vérifier le gain à fréquence nulle et les unités. Un pôle à partie réelle négative fait décroître sa contribution. Un terme intégrateur peut annuler l’écart statique à une consigne constante si la boucle reste stable, mais il modifie l’ordre et les pôles. Il faut donc contrôler la stabilité avant d’en déduire un gain de précision.

\section*{Fiche de révision}
Conditions initiales nulles pour un transfert, produit pour les blocs en série, somme signée pour le parallèle. En boucle négative, écrire $E=R-BY$ et $Y=GE$. Calculer séparément l’effet des perturbations, puis contrôler statiquement et dynamiquement le résultat.

\bigskip
{\footnotesize Adaptation française par Nexora CPGE. Source : Contributeurs de Wikibooks,
\href{https://en.wikibooks.org/wiki/Control_Systems/Feedback_Loops}{\textit{Control Systems}}.
\href{https://creativecommons.org/licenses/by-sa/4.0/}{CC BY-SA 4.0}.}
\end{document}
